% Template smoke test: IEEEtran conference with [unnumbered] option.
% Exercises the unnumbered fast path through the same heavy contexts
% as template-ieee-conference.tex (title, abstract, sections, short
% and long captions, figure*/table* wide floats, page marks). Adds a
% labeled note + \zebraref to verify the silent fallback to \ref.
\documentclass[conference]{IEEEtran}
\usepackage{lipsum}
\usepackage[unnumbered]{zebra}

\begin{document}
\title{Inline Notes (Unnumbered) Across IEEEtran\todo{title note}}
\author{\IEEEauthorblockN{A. Author}}
\maketitle

\begin{abstract}
\todo{abstract opener} A short abstract that contains an inline note
\todo{mid} and another at the end.\todo{abs end}
\end{abstract}

\begin{IEEEkeywords}
zebra, regression, IEEEtran\todo{kw}
\end{IEEEkeywords}

\section{Introduction\todo{intro section}}
\label{sec:intro}
Body paragraph with an inline \todo{\label{tnote:body}inline} note.
The fallback should resolve here: \zebraref{tnote:body}.
\lipsum[1][1-3]

\begin{figure}[t]
  \centering Placeholder
  \caption{Short caption\todo{short cap}.}
  \label{fig:short}
\end{figure}

\begin{figure*}[t]
  \centering Placeholder
  \caption{A long caption that exercises the \texttt{\textbackslash parbox}
    wrapping path inside \texttt{\textbackslash @makecaption}; the inline
    note must not corrupt the float queue.\todo{long cap}}
  \label{fig:wide}
\end{figure*}

\section{Tables}
\begin{table}[t]
  \centering
  \caption{Compact table caption with note\todo{tab cap}.}
  \label{tab:compact}
  \begin{tabular}{ll}\hline a & b \\ c & d \\\hline\end{tabular}
\end{table}

\begin{table*}[t]
  \centering
  \caption{A long table caption that wraps onto multiple lines so the
    parbox path is taken; the inline note must remain stable.\todo{wide tab cap}}
  \label{tab:wide}
  \begin{tabular}{lll}\hline x & y & z \\\hline\end{tabular}
\end{table*}

\lipsum[2-3]

% Cross-ref test: \zebraref must silently degrade to \ref when
% unnumbered. The label sec:intro is a regular section label.
See section~\zebraref{sec:intro} for context.
\end{document}
